% !TeX root = master_Manual.tex

%************************************************************************
% LaTeX Einstellungen (Simon Stephan 2016)
%
%
%************************************************************************


%************************************************************************
% Dokumenteinstellungen / Laden von Paketen
%
%************************************************************************

%\documentclass[%
%paper=A4,               % Papierformat
%twoside=false,          % beidseitiger Druck (Koma-Script)
%fontsize=11pt,          % Zeichensatzgroesse (Koma-Script)
%headings=normal,        % Ueberschriftengroesse (alternativ big oder small, Koma-Script)
%parskip=half+,          % Absatzkennzeichnung durch Zeilenabstand (Koma-Script)
%numbers=noenddot,       % letzte Ziffer der Nummerierung ohne nachfolgenden Punkt (Koma-Script)
%captions=tableheading,  % caption von Tabellen als Tabellenueberschrift formatieren (Koma-Script)
%bibliography=totoc,     % Literaturverzeichnis in Inhaltsverzeichnis (Koma-Script)
%listof=totoc,           % Tabellen- und Abbildungsverzeichnis in Inhaltsverzeichnis (Koma-Script)
%index=totoc             % Index in Inhaltsverzeichnis (Koma-Script)
%]
%{scrbook}               % Dokumentklasse (alternativ scrartcl oder scrreprt, Koma-Script) 

%Serifenlose Überschriften in KOMA-Script
\setkomafont{disposition}{\normalcolor\bfseries}

% Paket zur manuellen Randeinstellung
\usepackage{geometry}
\geometry{bindingoffset=10mm, left=15mm, right=15mm, top=20mm, bottom=25mm}

% Paket stellt die textblock-Umgebung zur Verfuegung, die frei auf der Seite
% platziert werden kann. Dabei koennen absolute Koordinaten (Papier) verwendet
% werden.
\usepackage[absolute]{textpos}


\usepackage{listings}
\usepackage{epsfig}
\usepackage{epstopdf}


% The bm pack­age de­fines a com­mand \bm which makes its ar­gu­ment bold
\usepackage{bm}


% Paket zur Gestaltung der Kolumnentitel
\usepackage{scrpage2}

% deutsche Trennungsregeln, deutsche Zusatzbefehle
\usepackage[english]{babel}

% Zeichenkodierung
\usepackage[utf8]{inputenc}

% deutsches Fontencoding, T1 Schriften
\usepackage[T1]{fontenc}

% Formatierung von Anführungszeichen --> \enquote Makro
% Option [german=guillemets] für französische Anführungszeichen
\usepackage{csquotes}

% Schriftfamilie ComputerModern Bright (serifenlos)
%\usepackage{cmbright}

% AMS-LaTeX Pakete fuer mathematischen Formelsatz. Die Option fleqn sorgt fuer
% linksbuendige Formeln, welche mit mathindent eingerueckt werden.
\usepackage[fleqn]{amsmath}
\usepackage{amssymb,amsthm}
\mathindent0.9cm

% stellt diverse Sonderzeichen zur Verfuegung z.B. \texteuro
\usepackage{textcomp}

% optischer Randausgleich (verschiedene Buchstaben, Satzzeichen, Bindestriche
% etc. werden leicht in den Seitenrand verschoben); automatische Skalierung von
% Buchstaben zur Verbesserung des Blocksatzes
\usepackage[tracking=true,spacing=true,kerning=true,final]{microtype}

% Silbentrennung wird auch bei linksbuendigem Text angewendet,
% z.B. im Index
\usepackage{ragged2e}

% das setspace package bietet die Moeglichkeit zur Aenderung des Durchschusses
% (Zeilenabstand im Blocksatz)
\usepackage{setspace}

% zusaetzliche Hinweise zur Seite auf der sich ein Verweis befindet
% z.B. mit dem Befehl \vref{label}
\usepackage[english]{varioref}

%% laedt die Bibliographie-Erweiterung biblatex
%\usepackage[ %
%backend=biber,         % Backend für biblatex [biber oder bibtex]
%style=authoryear-comp, % authoryear oder numeric (-comp --> compress-Option)
%maxbibnames=10,        % max. Anzahl Namen im Literaturverzeichnis 
%maxcitenames=1,        % max. Anzahl Namen bei authoryear Zitaten
%backref=true,          % Literaturverzeichnis mit Verweis Zitat-Seiten 
%%sorting=none,         % Verzeichnis nach Reihenfolge der Zitierung
%%dashed=false          % kein Strich im Verzeichnis bei Autor-Mehrfachnennung 
%]
%{biblatex}
%% Anpassung Einzug und Absatzabstand an Dokumentlayout
%\setlength{\bibhang}{0em} 
%\setlength{\bibitemsep}{\parskip}
%% Einbindung der bib-Datei
%\addbibresource{literatur.bib}

% Grafikunterstuetzung fuer LaTeX (Einfuegen von Bildern)
\usepackage{graphicx}

% Anordnung mehrerer Grafiken
\usepackage{subfig}

% Paket zum Drehen von Text, Tabellen, Grafiken etc.; das Paket stellt dafür
% die landscape-Umgebung zur Verfuegung
\usepackage{lscape}

% float-Objekte werden nur innerhalb der aktuellen section verschoben, ohne
% dieses Paket koennen floats auch in die naechste section uebernommen werden
\usepackage[section]{placeins}

% laedt das Grafikpaket tikz sowie einige Bibliotheken dazu
\usepackage{tikz}
\usetikzlibrary{arrows,shapes}


% Paket zum Setzen von Einheiten und zur Formatierung von Zahlen
% Einheiten werden immer "gerade" gesetzt. Zahlenformate, Trennzeichen etc.
% werden anhand der Spracheinstellungen (babel) ausgewaehlt (Manuelle
% Vorgaben sind ebenfalls moeglich --> siehe Paketdoku siunitx). 
% Befehle: \SI{Zahl}{Einheit}, \num{Zahl}
\usepackage[ngerman]{translator}
\usepackage[locale=DE, separate-uncertainty=true]{siunitx}

% Paket zum Erstellen neuer Spaltentypen
% --> Übergabe des Zelleninhaltes an Makros
\usepackage{collcell}

% mhchem Paket zum Setzen chemischer Formeln mit dem \ce Makro. Bei
% Verwendung im Text (ohne Mathe-Umgebung) wird die Schriftart automatisch
% an den umgebenden Text angepasst (Überschriften etc.).
\usepackage[version=3,arrows=pgf]{mhchem}
% Umbiegen des \ce Makros --> Für hyperref-Links wird der Quellcode der
% Formel als Linktext verwendet. Bei Nutzung des Makros in Überschriften 
% erscheinen sonst kryptische Linknamen, da im PDF dafür Sonderzeichen
% nicht erlaubt sind.
\let\oldce\ce
\renewcommand{\ce}[1]{\texorpdfstring{\oldce{#1}}{#1}}
% Neuer Spaltentyp für Tabellen mit \ce Makro
\newcolumntype{C}{>{\collectcell\ce}{l}<{\endcollectcell}}

% fuegt longtable-Umgebung ein (Tabellen mit moeglichem Seitenumbruch)
\usepackage{longtable}

% tabularx-Umgebung ermoeglicht Tabellen mit fester Breitenvorgabe,
% zusaetzlicher Spaltentyp X fuellt Tabelle auf Breitenvorgabe auf
\usepackage{tabularx}

% erweiterte Tabellenbefehle (zur Erstellung wirklich schoener Tabellen
% fuer den Buchdruck)
\usepackage{booktabs}

% erweiterte Moeglichkeiten zur Schriftformatierung in Tabellenspalten
\usepackage{array}

% Paket zum Zusammenfuegen von Tabellenzeilen
\usepackage{multirow}

% Paket zur Indexerstellung,  Indexeintraege mit \index{Name des Eintrages}
\usepackage{makeidx}
\makeindex

% erstellt Anweisungen, die nur fuer pdftex gelten
\usepackage{ifpdf}

% implementiert if Schleifen
\usepackage{ifthen}

% Paket zur automatischen Linkerstellung fuer Inhaltsverzeichnis, Literaturverweise etc.
% [plainpages=false] setzt getrennte Links fuer gleiche Seitenzahlen (z.B. fuer i und 1)
% [pdfpagelabels] geaendertes Anzeigeformat der Seitennummern im AcrobatReader (roemische
% Seitennummern werden angezeigt)
\usepackage[plainpages=false, pdfpagelabels]{hyperref}

% setzt Strafpunkte fuer Schusterjungen und Hurenkinder hoch 
\clubpenalty=10000
\widowpenalty=10000
\displaywidowpenalty=10000


%************************************************************************
% Anpassungen Kolumnentitel und Schriftarten
%
%************************************************************************

% Anpassung der Kolumnentitel, Seitennummerierung oben, Schriftgroessen
\clearscrheadfoot
\ohead{\headmark}
\ofoot[\pagemark]{\pagemark}
\pagestyle{scrheadings}
\addtokomafont{pageheadfoot}{\footnotesize}

% kleinere Schriftgroesse fuer Bildunterschriften/Tabellenueberschriften
\addtokomafont{caption}{\footnotesize}

% Anpassungen fuer das cmbright-package
% demibold (anstatt bold) als Fettschrift verwenden
%\renewcommand{\bfseries}{\fontseries{sb}\selectfont}
%% Befehl \textbx fuer Fette zur Verfuegung stellen 
%\newcommand{\textbx}[1]{{\fontseries{b}\selectfont#1}}
%\newcommand{\bxseries}{\fontseries{b}\selectfont}
% Ueberschriften in LatinModernSans-demicondensed
\newcommand{\headFont}{\fontfamily{lmss}\fontseries{sbc}\selectfont}
\addtokomafont{disposition}{\headFont}
% \mathsf fuer cmbright richtig definieren
% (sonst wird bei \mathsf CM-SansSerif verwendet)
\renewcommand{\mathsf}[1]{\mathrm{#1}}

\def\dcoeff{{\ooalign{\raisebox{.5mm}{\text{--}}\hspace{-1.9mm}}\hbox{\textit{\textbf{D}}}}}


%************************************************************************
% Definition Styles und Farbschemen
%
%************************************************************************

\makeatletter

% Voreinstellungen
\def\@docstyle{report}
\def\@color{white}

% Makro zur Auswahl von Style und Farbschema
\newcommand{\setlayout}[2]{%
  \def\@docstyle{#1}
  \def\@color{#2}
  \@setlayout
}


% Makro zur Umsetzung von Style und Farbschema
\newcommand{\@setlayout}{%
  % Farbschema weiss
  \ifthenelse{\equal{\@color}{white}}{% 
  }
}


%************************************************************************
% Titelseite
%
%************************************************************************

% Variablen fuer Benutzereinstellungen
\def\@faculty{}
\def\@chair{}
\def\@subTitle{}
\def\@titleFoot{}

\newcommand{\faculty}[1]{\def\@faculty{#1}}
\newcommand{\chair}[1]{\def\@chair{, #1}}
\newcommand{\subTitle}[1]{\def\@subTitle{#1}}
\newcommand{\titleFoot}[1]{\def\@titleFoot{#1}}

\renewcommand{\maketitle}{%
  \begin{titlepage}
    \null%
    \begin{textblock*}{297mm}(0mm,0mm)%
    {\setlength{\unitlength}{1mm}%
    \noindent%
    \begin{picture}(0,0)%     

      \ifthenelse{\equal{\@color}{white}}{%  

        \def\@headfontcolor{\color{black}}
        \def\@textfontcolor{\color{black}}
        \def\@logofile{grafiken/logos/Logo_LTD_Text}
      }

      
      % Minipage als Bereich fuer den Seiteninhalt     
      \put(30,-44){\@textfontcolor
        \begin{minipage}[t][140mm][l]{160mm}

            \vfill
            \vfill
            \vfill
            \vskip 2em
            \begin{center}
              {\headFont\Huge\@title}
              \vskip 1em
              \textbf{\Large\@subtitle}
            
            \vskip 1.5em
              \Large\@subTitle
            
            \vskip 3em
            {\large\@date}      
            \vfill
            \phantom{x}
            \end{center}

       \end{minipage} 
      }

      % Infoblock Fusszeile (in fester Schriftgroesse)
      \put(30,-280){\@headfontcolor\fontsize{8pt}{10pt}\selectfont
        \begin{minipage}[t][21mm][l]{160mm}
          {\setlength{\tabcolsep}{1.5em} %groesserer Spaltenabstand
          \@titleFoot
          }
        \end{minipage}
      } 
    \end{picture}
  }
  \end{textblock*}
\end{titlepage}
}


%************************************************************************
% pdf Dokumenteninformationen (Datenübernahme aus Titelangaben)
%
%************************************************************************

\newcommand{\makepdfinfos}{%
  \hypersetup{  
    pdfauthor = {\@author},
    pdftitle = {\@title},
    pdfsubject = {\@subTitle},   
    pdfkeywords = {},
    bookmarksnumbered = {true}, % Lesezeichen im pdf mit Nummerierung
    pdfpagemode = {UseOutlines}, % UseOutlines, UseThumbs, None
    pdfstartview = {FitH}, % FitH, FitV oder XYZ null null 0.9
    pdffitwindow = {true}, % Fenstergröße anpassen
  }
}



\makeatother

%%% Local Variables: 
%%% mode: latex
%%% TeX-master: "master"
%%% coding: utf-8
%%% End:
